% ═══════════════════════════════════════════════════════════════════════════════
%  Stile TikZ che riproduce il linguaggio visivo della figura-framework MLISE.
%  Ogni valore viene da una misura su ClusteringFramework.png — vedi STYLE.md.
%
%  Alternativa vettoriale al PNG da 6,36 MB: il testo usa i font del documento,
%  la figura resta editabile e pesa qualche kB.
%
%  Uso nel preambolo (IEEEtran conference, come il tuo paper.tex):
%      \usepackage{tikz}
%      \usetikzlibrary{arrows.meta, shapes.geometric, positioning, fit, backgrounds}
%      % ═══════════════════════════════════════════════════════════════════════════════
%  Stile TikZ che riproduce il linguaggio visivo della figura-framework MLISE.
%  Ogni valore viene da una misura su ClusteringFramework.png — vedi STYLE.md.
%
%  Alternativa vettoriale al PNG da 6,36 MB: il testo usa i font del documento,
%  la figura resta editabile e pesa qualche kB.
%
%  Uso nel preambolo (IEEEtran conference, come il tuo paper.tex):
%      \usepackage{tikz}
%      \usetikzlibrary{arrows.meta, shapes.geometric, positioning, fit, backgrounds}
%      % ═══════════════════════════════════════════════════════════════════════════════
%  Stile TikZ che riproduce il linguaggio visivo della figura-framework MLISE.
%  Ogni valore viene da una misura su ClusteringFramework.png — vedi STYLE.md.
%
%  Alternativa vettoriale al PNG da 6,36 MB: il testo usa i font del documento,
%  la figura resta editabile e pesa qualche kB.
%
%  Uso nel preambolo (IEEEtran conference, come il tuo paper.tex):
%      \usepackage{tikz}
%      \usetikzlibrary{arrows.meta, shapes.geometric, positioning, fit, backgrounds}
%      % ═══════════════════════════════════════════════════════════════════════════════
%  Stile TikZ che riproduce il linguaggio visivo della figura-framework MLISE.
%  Ogni valore viene da una misura su ClusteringFramework.png — vedi STYLE.md.
%
%  Alternativa vettoriale al PNG da 6,36 MB: il testo usa i font del documento,
%  la figura resta editabile e pesa qualche kB.
%
%  Uso nel preambolo (IEEEtran conference, come il tuo paper.tex):
%      \usepackage{tikz}
%      \usetikzlibrary{arrows.meta, shapes.geometric, positioning, fit, backgrounds}
%      \input{framework_style}
%
%  ⚠️ NON COMPILATO: sulla macchina dove è stato scritto non c'è pdflatex.
%     Compilalo una volta prima di fidartene.
% ═══════════════════════════════════════════════════════════════════════════════

% ── 1. palette: colori PURI, l'effetto pastello lo fa l'opacità ────────────────
\definecolor{fwStageA}{HTML}{FF8029}   % sfondo stage 1   (opacità 0.19)
\definecolor{fwStageB}{HTML}{FFFF00}   % sfondo stage 2   (opacità 0.19)
\definecolor{fwPanel} {HTML}{71CEDC}   % pannello         (opacità 0.59)
\definecolor{fwPanelB}{HTML}{71E2D5}   % box decoder      (opacità 0.59)
\definecolor{fwEnc}   {HTML}{FF9A9A}   % encoder          (opacità 0.49)
\definecolor{fwInput} {HTML}{AEC6FF}   % ingresso         (opacità 0.49)
\definecolor{fwView}  {HTML}{90EF90}   % vista            (opacità 0.49)
\definecolor{fwInner} {HTML}{F9F9FA}   % interno ellissi
% accenti pieni per le illustrazioni interne
\definecolor{fwBlue}  {HTML}{4567A9}
\definecolor{fwOrange}{HTML}{D17920}
\definecolor{fwGreen} {HTML}{2B7B4F}
\definecolor{fwPurple}{HTML}{7B3A99}
\definecolor{fwRed}   {HTML}{B0302F}

% ── 2. tipografia: pesante, corsiva, con grazie ───────────────────────────────
% pbk = URW Bookman nella nomenclatura PostScript standard. Cambia il font SOLO
% dentro la figura, invece di \usepackage{bookman} che lo cambierebbe ovunque.
% Se pbk non è installato (pacchetto fonts-urw-base35 / tex-gyre), togli
% \fontfamily{pbk} e resta il corsivo grassetto del documento.
\newcommand{\fwFont}{\fontfamily{pbk}\selectfont\bfseries\itshape}
\newcommand{\fwText}[1]{{\fwFont #1}}

\tikzset{
  % ── 3. tratto UNIFORME: ~1 pt alla dimensione di resa, su ogni forma ────────
  fw/.style        = {line width=1pt, draw=black, line join=round, line cap=round},
  % ── 4. forme ───────────────────────────────────────────────────────────────
  % raggio d'angolo = 35% dell'altezza del box: con altezza 8mm -> 2.8mm
  fw box/.style    = {fw, rounded corners=2.8mm, inner sep=2.2mm,
                      align=center, font=\footnotesize, minimum height=8mm},
  fw input/.style  = {fw box, fill=fwInput,  fill opacity=0.49, text opacity=1},
  fw view/.style   = {fw box, fill=fwView,   fill opacity=0.49, text opacity=1},
  fw panel/.style  = {fw box, fill=fwPanel,  fill opacity=0.59, text opacity=1},
  fw decoder/.style= {fw box, fill=fwPanelB, fill opacity=0.59, text opacity=1},
  % encoder: trapezio rastremato verso il basso
  fw encoder/.style= {fw, trapezium, trapezium left angle=72, trapezium right angle=72,
                      trapezium stretches body, shape border rotate=180,
                      fill=fwEnc, fill opacity=0.49, text opacity=1,
                      align=center, font=\footnotesize, inner sep=2mm, minimum height=9mm},
  % contenitore delle illustrazioni concettuali
  fw bubble/.style = {fw, ellipse, fill=fwInner, fill opacity=0.99, text opacity=1,
                      align=center, font=\scriptsize, inner sep=2.4mm},
  % ── 5. sfondi di stage: grandi, arrotondati, SENZA contorno ────────────────
  fw stage a/.style= {rounded corners=3mm, fill=fwStageA, fill opacity=0.19, draw=none},
  fw stage b/.style= {rounded corners=3mm, fill=fwStageB, fill opacity=0.19, draw=none},
  fw stage tag/.style = {font=\small, anchor=south west, inner sep=1mm},
  % ── 6. frecce: punta triangolare barbata ───────────────────────────────────
  fw arrow/.style  = {fw, -{Stealth[length=2.6mm, width=2.2mm, inset=0.7mm]}},
  fw dash/.style   = {fw, dash pattern=on 1.4mm off 1mm},
  % nodo di loss: terminale, grande
  fw loss/.style   = {font=\normalsize, inner sep=0.6mm},
}

% Etichetta di loss nello stesso peso della figura: \fwLoss{cross} -> L_cross
\newcommand{\fwLoss}[1]{$\boldsymbol{\mathcal{L}}_{\text{\scriptsize\itshape #1}}$}

% ═══════════════════════════════════════════════════════════════════════════════
%  Scheletro d'uso — la composizione dell'originale: ingresso in alto, ventaglio
%  su N viste, una riga di encoder, le rappresentazioni, poi DUE regioni di stage
%  affiancate ciascuna con i propri termini di loss che convergono su una finale.
%
%  \begin{figure}[t]\centering
%  \begin{tikzpicture}[node distance=6mm and 8mm]
%    \node[fw input] (aug) {\fwText{Stochastic Augmentation}\\\fwText{\& Random Masking}};
%    \node[fw view, below left=9mm and 16mm of aug]  (v1) {\fwText{View 1}};
%    \node[fw view, below=9mm of aug]                (v2) {\fwText{View 2}};
%    \node[fw view, below right=9mm and 16mm of aug] (vn) {\fwText{View N}};
%    \node[fw encoder, below=6mm of v1] (e1) {\fwText{Mamba}\\\fwText{Encoder}};
%    \foreach \s/\t in {aug/v1, aug/v2, aug/vn, v1/e1} \draw[fw arrow] (\s) -- (\t);
%    \node[fw loss] (lp) {\fwLoss{pre}};
%    % lo sfondo di stage si disegna DOPO, sul layer background, con `fit`
%    \begin{scope}[on background layer]
%      \node[fw stage a, fit=(e1)(lp), inner sep=3mm] (sa) {};
%    \end{scope}
%    \node[fw stage tag, below right] at (sa.south west) {\fwText{Pretraining Stage}};
%  \end{tikzpicture}
%  \caption{...}\end{figure}
% ═══════════════════════════════════════════════════════════════════════════════

%
%  ⚠️ NON COMPILATO: sulla macchina dove è stato scritto non c'è pdflatex.
%     Compilalo una volta prima di fidartene.
% ═══════════════════════════════════════════════════════════════════════════════

% ── 1. palette: colori PURI, l'effetto pastello lo fa l'opacità ────────────────
\definecolor{fwStageA}{HTML}{FF8029}   % sfondo stage 1   (opacità 0.19)
\definecolor{fwStageB}{HTML}{FFFF00}   % sfondo stage 2   (opacità 0.19)
\definecolor{fwPanel} {HTML}{71CEDC}   % pannello         (opacità 0.59)
\definecolor{fwPanelB}{HTML}{71E2D5}   % box decoder      (opacità 0.59)
\definecolor{fwEnc}   {HTML}{FF9A9A}   % encoder          (opacità 0.49)
\definecolor{fwInput} {HTML}{AEC6FF}   % ingresso         (opacità 0.49)
\definecolor{fwView}  {HTML}{90EF90}   % vista            (opacità 0.49)
\definecolor{fwInner} {HTML}{F9F9FA}   % interno ellissi
% accenti pieni per le illustrazioni interne
\definecolor{fwBlue}  {HTML}{4567A9}
\definecolor{fwOrange}{HTML}{D17920}
\definecolor{fwGreen} {HTML}{2B7B4F}
\definecolor{fwPurple}{HTML}{7B3A99}
\definecolor{fwRed}   {HTML}{B0302F}

% ── 2. tipografia: pesante, corsiva, con grazie ───────────────────────────────
% pbk = URW Bookman nella nomenclatura PostScript standard. Cambia il font SOLO
% dentro la figura, invece di \usepackage{bookman} che lo cambierebbe ovunque.
% Se pbk non è installato (pacchetto fonts-urw-base35 / tex-gyre), togli
% \fontfamily{pbk} e resta il corsivo grassetto del documento.
\newcommand{\fwFont}{\fontfamily{pbk}\selectfont\bfseries\itshape}
\newcommand{\fwText}[1]{{\fwFont #1}}

\tikzset{
  % ── 3. tratto UNIFORME: ~1 pt alla dimensione di resa, su ogni forma ────────
  fw/.style        = {line width=1pt, draw=black, line join=round, line cap=round},
  % ── 4. forme ───────────────────────────────────────────────────────────────
  % raggio d'angolo = 35% dell'altezza del box: con altezza 8mm -> 2.8mm
  fw box/.style    = {fw, rounded corners=2.8mm, inner sep=2.2mm,
                      align=center, font=\footnotesize, minimum height=8mm},
  fw input/.style  = {fw box, fill=fwInput,  fill opacity=0.49, text opacity=1},
  fw view/.style   = {fw box, fill=fwView,   fill opacity=0.49, text opacity=1},
  fw panel/.style  = {fw box, fill=fwPanel,  fill opacity=0.59, text opacity=1},
  fw decoder/.style= {fw box, fill=fwPanelB, fill opacity=0.59, text opacity=1},
  % encoder: trapezio rastremato verso il basso
  fw encoder/.style= {fw, trapezium, trapezium left angle=72, trapezium right angle=72,
                      trapezium stretches body, shape border rotate=180,
                      fill=fwEnc, fill opacity=0.49, text opacity=1,
                      align=center, font=\footnotesize, inner sep=2mm, minimum height=9mm},
  % contenitore delle illustrazioni concettuali
  fw bubble/.style = {fw, ellipse, fill=fwInner, fill opacity=0.99, text opacity=1,
                      align=center, font=\scriptsize, inner sep=2.4mm},
  % ── 5. sfondi di stage: grandi, arrotondati, SENZA contorno ────────────────
  fw stage a/.style= {rounded corners=3mm, fill=fwStageA, fill opacity=0.19, draw=none},
  fw stage b/.style= {rounded corners=3mm, fill=fwStageB, fill opacity=0.19, draw=none},
  fw stage tag/.style = {font=\small, anchor=south west, inner sep=1mm},
  % ── 6. frecce: punta triangolare barbata ───────────────────────────────────
  fw arrow/.style  = {fw, -{Stealth[length=2.6mm, width=2.2mm, inset=0.7mm]}},
  fw dash/.style   = {fw, dash pattern=on 1.4mm off 1mm},
  % nodo di loss: terminale, grande
  fw loss/.style   = {font=\normalsize, inner sep=0.6mm},
}

% Etichetta di loss nello stesso peso della figura: \fwLoss{cross} -> L_cross
\newcommand{\fwLoss}[1]{$\boldsymbol{\mathcal{L}}_{\text{\scriptsize\itshape #1}}$}

% ═══════════════════════════════════════════════════════════════════════════════
%  Scheletro d'uso — la composizione dell'originale: ingresso in alto, ventaglio
%  su N viste, una riga di encoder, le rappresentazioni, poi DUE regioni di stage
%  affiancate ciascuna con i propri termini di loss che convergono su una finale.
%
%  \begin{figure}[t]\centering
%  \begin{tikzpicture}[node distance=6mm and 8mm]
%    \node[fw input] (aug) {\fwText{Stochastic Augmentation}\\\fwText{\& Random Masking}};
%    \node[fw view, below left=9mm and 16mm of aug]  (v1) {\fwText{View 1}};
%    \node[fw view, below=9mm of aug]                (v2) {\fwText{View 2}};
%    \node[fw view, below right=9mm and 16mm of aug] (vn) {\fwText{View N}};
%    \node[fw encoder, below=6mm of v1] (e1) {\fwText{Mamba}\\\fwText{Encoder}};
%    \foreach \s/\t in {aug/v1, aug/v2, aug/vn, v1/e1} \draw[fw arrow] (\s) -- (\t);
%    \node[fw loss] (lp) {\fwLoss{pre}};
%    % lo sfondo di stage si disegna DOPO, sul layer background, con `fit`
%    \begin{scope}[on background layer]
%      \node[fw stage a, fit=(e1)(lp), inner sep=3mm] (sa) {};
%    \end{scope}
%    \node[fw stage tag, below right] at (sa.south west) {\fwText{Pretraining Stage}};
%  \end{tikzpicture}
%  \caption{...}\end{figure}
% ═══════════════════════════════════════════════════════════════════════════════

%
%  ⚠️ NON COMPILATO: sulla macchina dove è stato scritto non c'è pdflatex.
%     Compilalo una volta prima di fidartene.
% ═══════════════════════════════════════════════════════════════════════════════

% ── 1. palette: colori PURI, l'effetto pastello lo fa l'opacità ────────────────
\definecolor{fwStageA}{HTML}{FF8029}   % sfondo stage 1   (opacità 0.19)
\definecolor{fwStageB}{HTML}{FFFF00}   % sfondo stage 2   (opacità 0.19)
\definecolor{fwPanel} {HTML}{71CEDC}   % pannello         (opacità 0.59)
\definecolor{fwPanelB}{HTML}{71E2D5}   % box decoder      (opacità 0.59)
\definecolor{fwEnc}   {HTML}{FF9A9A}   % encoder          (opacità 0.49)
\definecolor{fwInput} {HTML}{AEC6FF}   % ingresso         (opacità 0.49)
\definecolor{fwView}  {HTML}{90EF90}   % vista            (opacità 0.49)
\definecolor{fwInner} {HTML}{F9F9FA}   % interno ellissi
% accenti pieni per le illustrazioni interne
\definecolor{fwBlue}  {HTML}{4567A9}
\definecolor{fwOrange}{HTML}{D17920}
\definecolor{fwGreen} {HTML}{2B7B4F}
\definecolor{fwPurple}{HTML}{7B3A99}
\definecolor{fwRed}   {HTML}{B0302F}

% ── 2. tipografia: pesante, corsiva, con grazie ───────────────────────────────
% pbk = URW Bookman nella nomenclatura PostScript standard. Cambia il font SOLO
% dentro la figura, invece di \usepackage{bookman} che lo cambierebbe ovunque.
% Se pbk non è installato (pacchetto fonts-urw-base35 / tex-gyre), togli
% \fontfamily{pbk} e resta il corsivo grassetto del documento.
\newcommand{\fwFont}{\fontfamily{pbk}\selectfont\bfseries\itshape}
\newcommand{\fwText}[1]{{\fwFont #1}}

\tikzset{
  % ── 3. tratto UNIFORME: ~1 pt alla dimensione di resa, su ogni forma ────────
  fw/.style        = {line width=1pt, draw=black, line join=round, line cap=round},
  % ── 4. forme ───────────────────────────────────────────────────────────────
  % raggio d'angolo = 35% dell'altezza del box: con altezza 8mm -> 2.8mm
  fw box/.style    = {fw, rounded corners=2.8mm, inner sep=2.2mm,
                      align=center, font=\footnotesize, minimum height=8mm},
  fw input/.style  = {fw box, fill=fwInput,  fill opacity=0.49, text opacity=1},
  fw view/.style   = {fw box, fill=fwView,   fill opacity=0.49, text opacity=1},
  fw panel/.style  = {fw box, fill=fwPanel,  fill opacity=0.59, text opacity=1},
  fw decoder/.style= {fw box, fill=fwPanelB, fill opacity=0.59, text opacity=1},
  % encoder: trapezio rastremato verso il basso
  fw encoder/.style= {fw, trapezium, trapezium left angle=72, trapezium right angle=72,
                      trapezium stretches body, shape border rotate=180,
                      fill=fwEnc, fill opacity=0.49, text opacity=1,
                      align=center, font=\footnotesize, inner sep=2mm, minimum height=9mm},
  % contenitore delle illustrazioni concettuali
  fw bubble/.style = {fw, ellipse, fill=fwInner, fill opacity=0.99, text opacity=1,
                      align=center, font=\scriptsize, inner sep=2.4mm},
  % ── 5. sfondi di stage: grandi, arrotondati, SENZA contorno ────────────────
  fw stage a/.style= {rounded corners=3mm, fill=fwStageA, fill opacity=0.19, draw=none},
  fw stage b/.style= {rounded corners=3mm, fill=fwStageB, fill opacity=0.19, draw=none},
  fw stage tag/.style = {font=\small, anchor=south west, inner sep=1mm},
  % ── 6. frecce: punta triangolare barbata ───────────────────────────────────
  fw arrow/.style  = {fw, -{Stealth[length=2.6mm, width=2.2mm, inset=0.7mm]}},
  fw dash/.style   = {fw, dash pattern=on 1.4mm off 1mm},
  % nodo di loss: terminale, grande
  fw loss/.style   = {font=\normalsize, inner sep=0.6mm},
}

% Etichetta di loss nello stesso peso della figura: \fwLoss{cross} -> L_cross
\newcommand{\fwLoss}[1]{$\boldsymbol{\mathcal{L}}_{\text{\scriptsize\itshape #1}}$}

% ═══════════════════════════════════════════════════════════════════════════════
%  Scheletro d'uso — la composizione dell'originale: ingresso in alto, ventaglio
%  su N viste, una riga di encoder, le rappresentazioni, poi DUE regioni di stage
%  affiancate ciascuna con i propri termini di loss che convergono su una finale.
%
%  \begin{figure}[t]\centering
%  \begin{tikzpicture}[node distance=6mm and 8mm]
%    \node[fw input] (aug) {\fwText{Stochastic Augmentation}\\\fwText{\& Random Masking}};
%    \node[fw view, below left=9mm and 16mm of aug]  (v1) {\fwText{View 1}};
%    \node[fw view, below=9mm of aug]                (v2) {\fwText{View 2}};
%    \node[fw view, below right=9mm and 16mm of aug] (vn) {\fwText{View N}};
%    \node[fw encoder, below=6mm of v1] (e1) {\fwText{Mamba}\\\fwText{Encoder}};
%    \foreach \s/\t in {aug/v1, aug/v2, aug/vn, v1/e1} \draw[fw arrow] (\s) -- (\t);
%    \node[fw loss] (lp) {\fwLoss{pre}};
%    % lo sfondo di stage si disegna DOPO, sul layer background, con `fit`
%    \begin{scope}[on background layer]
%      \node[fw stage a, fit=(e1)(lp), inner sep=3mm] (sa) {};
%    \end{scope}
%    \node[fw stage tag, below right] at (sa.south west) {\fwText{Pretraining Stage}};
%  \end{tikzpicture}
%  \caption{...}\end{figure}
% ═══════════════════════════════════════════════════════════════════════════════

%
%  ⚠️ NON COMPILATO: sulla macchina dove è stato scritto non c'è pdflatex.
%     Compilalo una volta prima di fidartene.
% ═══════════════════════════════════════════════════════════════════════════════

% ── 1. palette: colori PURI, l'effetto pastello lo fa l'opacità ────────────────
\definecolor{fwStageA}{HTML}{FF8029}   % sfondo stage 1   (opacità 0.19)
\definecolor{fwStageB}{HTML}{FFFF00}   % sfondo stage 2   (opacità 0.19)
\definecolor{fwPanel} {HTML}{71CEDC}   % pannello         (opacità 0.59)
\definecolor{fwPanelB}{HTML}{71E2D5}   % box decoder      (opacità 0.59)
\definecolor{fwEnc}   {HTML}{FF9A9A}   % encoder          (opacità 0.49)
\definecolor{fwInput} {HTML}{AEC6FF}   % ingresso         (opacità 0.49)
\definecolor{fwView}  {HTML}{90EF90}   % vista            (opacità 0.49)
\definecolor{fwInner} {HTML}{F9F9FA}   % interno ellissi
% accenti pieni per le illustrazioni interne
\definecolor{fwBlue}  {HTML}{4567A9}
\definecolor{fwOrange}{HTML}{D17920}
\definecolor{fwGreen} {HTML}{2B7B4F}
\definecolor{fwPurple}{HTML}{7B3A99}
\definecolor{fwRed}   {HTML}{B0302F}

% ── 2. tipografia: pesante, corsiva, con grazie ───────────────────────────────
% pbk = URW Bookman nella nomenclatura PostScript standard. Cambia il font SOLO
% dentro la figura, invece di \usepackage{bookman} che lo cambierebbe ovunque.
% Se pbk non è installato (pacchetto fonts-urw-base35 / tex-gyre), togli
% \fontfamily{pbk} e resta il corsivo grassetto del documento.
\newcommand{\fwFont}{\fontfamily{pbk}\selectfont\bfseries\itshape}
\newcommand{\fwText}[1]{{\fwFont #1}}

\tikzset{
  % ── 3. tratto UNIFORME: ~1 pt alla dimensione di resa, su ogni forma ────────
  fw/.style        = {line width=1pt, draw=black, line join=round, line cap=round},
  % ── 4. forme ───────────────────────────────────────────────────────────────
  % raggio d'angolo = 35% dell'altezza del box: con altezza 8mm -> 2.8mm
  fw box/.style    = {fw, rounded corners=2.8mm, inner sep=2.2mm,
                      align=center, font=\footnotesize, minimum height=8mm},
  fw input/.style  = {fw box, fill=fwInput,  fill opacity=0.49, text opacity=1},
  fw view/.style   = {fw box, fill=fwView,   fill opacity=0.49, text opacity=1},
  fw panel/.style  = {fw box, fill=fwPanel,  fill opacity=0.59, text opacity=1},
  fw decoder/.style= {fw box, fill=fwPanelB, fill opacity=0.59, text opacity=1},
  % encoder: trapezio rastremato verso il basso
  fw encoder/.style= {fw, trapezium, trapezium left angle=72, trapezium right angle=72,
                      trapezium stretches body, shape border rotate=180,
                      fill=fwEnc, fill opacity=0.49, text opacity=1,
                      align=center, font=\footnotesize, inner sep=2mm, minimum height=9mm},
  % contenitore delle illustrazioni concettuali
  fw bubble/.style = {fw, ellipse, fill=fwInner, fill opacity=0.99, text opacity=1,
                      align=center, font=\scriptsize, inner sep=2.4mm},
  % ── 5. sfondi di stage: grandi, arrotondati, SENZA contorno ────────────────
  fw stage a/.style= {rounded corners=3mm, fill=fwStageA, fill opacity=0.19, draw=none},
  fw stage b/.style= {rounded corners=3mm, fill=fwStageB, fill opacity=0.19, draw=none},
  fw stage tag/.style = {font=\small, anchor=south west, inner sep=1mm},
  % ── 6. frecce: punta triangolare barbata ───────────────────────────────────
  fw arrow/.style  = {fw, -{Stealth[length=2.6mm, width=2.2mm, inset=0.7mm]}},
  fw dash/.style   = {fw, dash pattern=on 1.4mm off 1mm},
  % nodo di loss: terminale, grande
  fw loss/.style   = {font=\normalsize, inner sep=0.6mm},
}

% Etichetta di loss nello stesso peso della figura: \fwLoss{cross} -> L_cross
\newcommand{\fwLoss}[1]{$\boldsymbol{\mathcal{L}}_{\text{\scriptsize\itshape #1}}$}

% ═══════════════════════════════════════════════════════════════════════════════
%  Scheletro d'uso — la composizione dell'originale: ingresso in alto, ventaglio
%  su N viste, una riga di encoder, le rappresentazioni, poi DUE regioni di stage
%  affiancate ciascuna con i propri termini di loss che convergono su una finale.
%
%  \begin{figure}[t]\centering
%  \begin{tikzpicture}[node distance=6mm and 8mm]
%    \node[fw input] (aug) {\fwText{Stochastic Augmentation}\\\fwText{\& Random Masking}};
%    \node[fw view, below left=9mm and 16mm of aug]  (v1) {\fwText{View 1}};
%    \node[fw view, below=9mm of aug]                (v2) {\fwText{View 2}};
%    \node[fw view, below right=9mm and 16mm of aug] (vn) {\fwText{View N}};
%    \node[fw encoder, below=6mm of v1] (e1) {\fwText{Mamba}\\\fwText{Encoder}};
%    \foreach \s/\t in {aug/v1, aug/v2, aug/vn, v1/e1} \draw[fw arrow] (\s) -- (\t);
%    \node[fw loss] (lp) {\fwLoss{pre}};
%    % lo sfondo di stage si disegna DOPO, sul layer background, con `fit`
%    \begin{scope}[on background layer]
%      \node[fw stage a, fit=(e1)(lp), inner sep=3mm] (sa) {};
%    \end{scope}
%    \node[fw stage tag, below right] at (sa.south west) {\fwText{Pretraining Stage}};
%  \end{tikzpicture}
%  \caption{...}\end{figure}
% ═══════════════════════════════════════════════════════════════════════════════
